\section{Morphisme d’incidence hexade--octade et conservation des canaux orientés \(90/120\)}

La dodécade et l’alignement du
\cref{sec:w12-w24-elevacion} étant fixés, soit \(H\) une hexade dans les coordonnées HMT,
soit \(\ell(H)\subset D\) son image et soit
\(O_H=\iota_D(\ell(H))\) son unique relèvement. En marquant un point
et une paire de points de l’hexade résiduelle, on obtient
\[
F_6(H)=\ell(H)\times\binom{\ell(H)}2,
\qquad
|F_6|=6\binom62=90.
\]
Si le point marqué peut parcourir l’octade :
\[
F_8(H)=O_H\times\binom{\ell(H)}2,
\qquad
|F_8|=8\binom62=120.
\]
L’inclusion \(\ell(H)\hookrightarrow O_H\) induit le morphisme d’incidence
\[
 \mathcal I_{6\to8}:F_6(H)\hookrightarrow F_8(H).
\]
En normalisant la différence par le facteur
choisi \(24\binom62\), on obtient
\[
\boxed{
\frac{90-120}{24\binom62}
=\frac{6-8}{24}
=-\frac1{12}.}
\]
L’égalité est un défaut normalisé de l’interface combinatoire
\(6\to8\). Le facteur \(24\binom62\) fait partie de la définition de cette
normalisation ; l’inclusion à elle seule détermine la différence \(-30\).
L’opérateur \(\mathcal I_{6\to8}\) a pour domaine les configurations
marquées d’incidence ; ce n’est ni l’extracteur transversal de l’état
dodécaphasique ni la famille des moments angulaires.
